\documentclass[10pt,a4paper]{amsart}
\usepackage[margin=19mm]{geometry}
\usepackage{amsmath,amssymb,mathtools}
\usepackage[scr=rsfs,cal=euler]{mathalfa}
\usepackage{xfrac}
\usepackage[unicode,hidelinks,bookmarks=false]{hyperref}
\newcommand{\ie}{i.e.\ }
\newcommand{\cf}{cf.\ }
\newcommand{\resp}{respectively\ }
\newcommand{\Subsection}{\subsection}
\providecommand{\nobreakdash}{\nobreak\hbox{-}}
\setlength{\emergencystretch}{2em}
\multlinegap0pt
\usepackage{bm}
\usepackage{bbm}
\usepackage{dsfont}
\usepackage{xspace}
\usepackage{cancel}
\usepackage{mathtools}
\usepackage{amsthm}
\makeatletter
\@ifundefined{proposition}{\newtheorem{proposition}{Proposition}}{}
\makeatother
\title[Learning collective variables that preserve transition rates]{Learning collective variables that preserve transition rates}
\author{Shashank Sule, Arnav Mehta, Maria K. Cameron}
\date{Source: arXiv:2506.01222v2 (CC BY 4.0)}
\begin{document}
\maketitle

\noindent\emph{Source and licence.} Learning collective variables that preserve transition rates. Version arXiv:2506.01222v2 (CC BY 4.0), distributed under the Creative Commons Attribution 4.0 International licence, as recorded in the arXiv licence metadata for this exact version and shown on the abstract page https://arxiv.org/abs/2506.01222v2. Full rights notice: licenses/arxiv-2506.01222-CC-BY-4.0.txt.\par\smallskip

\noindent\textit{Editorial scope.} Counted unit: the Proposition whose source label is \texttt{prop: oc poc equivalence} in the article (source0.tex lines 530--533, source label \texttt{prop: oc poc equivalence}), delivered together with the article's own complete original proof of it (source0.tex lines 535--579, one \texttt{proof} environment of 45 lines). No mathematics is written, repaired or re-derived: statement and proof are the article's own text line for line. Required context retained uncounted: eq: split potential (lines 91--93); eq: orthogonality condition (lines 101--103); eq: orthogonality condition 2 (lines 381--383). Every definition, display and label of those lines is retained unchanged. Omitted from this package and not redistributed: 1--90, 94--100, 104--380, 384--529, 534--534, 580--1333. Nothing inside a retained excerpt is omitted or altered: every retained body is delivered exactly as the article's lines are written, so no omission receipt of any kind accompanies this package. Citations: the retained excerpts carry no citation command at all, so no bibliography entry of the article is transcribed and no reference is redistributed. Reference adaptation: none was needed and none was made; every reference inside the delivered unit resolves inside it, so no unresolved reference or citation is produced. Printed slips kept literal: no printed word, symbol, hypothesis or number of the counted statement or of its proof was altered. Layout and class: the article's own class is \texttt{article} with packages including graphicx, authblk, amsthm, amsfonts, amssymb, amsmath, booktabs, multirow, multicol, longtable, biblatex, geometry, algorithm2e, hyperref; the wrapper is the lane's pinned \texttt{amsart} editorial wrapper at 10pt on A4, which restates the theorem-like environments the retained text needs and adds the article's own macro definitions that the retained text needs (none), with extra wrapper packages (bm, bbm, dsfont, xspace, cancel, mathtools). The delivered numbering is the unit's own. Assets: no retained span contains a figure environment, a graphics-inclusion command, a PDF-page inclusion, a table or a TikZ picture; no image file is copied into this package, the package declares no asset dependency and no image file is redistributed with it. Licence: version arXiv:2506.01222v2 of this article is distributed under the Creative Commons Attribution 4.0 International licence, as recorded in the arXiv licence metadata for this exact version and shown on the abstract page https://arxiv.org/abs/2506.01222v2. A full rights notice is delivered at licenses/arxiv-2506.01222-CC-BY-4.0.txt. Characters: every retained excerpt is pure ASCII, so the delivered engine, XeLaTeX with its default Latin Modern text font, reports no missing character.
\begin{align}
    V = V_0 + \epsilon^{-1}V_1. \label{eq: split potential} 
\end{align}

\begin{align}
    D\xi\nabla V_1 = 0. \label{eq: orthogonality condition} \tag{OC}
\end{align}

\begin{equation}
    (I - \Pi)^\top \nabla V_1 \equiv 0. \label{eq: orthogonality condition 2} \tag{POC}
\end{equation}

\begin{proposition}
\label{prop: oc poc equivalence}
    Let $\xi$ be a collective variable and $V$ be given by the decomposition $V= V_0 + \epsilon^{-1}V_{1}$ as in \eqref{eq: split potential}. Then \eqref{eq: orthogonality condition} and \eqref{eq: orthogonality condition 2} are equivalent. 
\end{proposition}

\begin{proof}
%    We show that \eqref{eq: orthogonality condition} $\iff$ \eqref{eq: orthogonality condition 2} for every $x \in \mathbb{R}^{N}$. %To that end, first to shorten notation we set
    Let $J := D\xi(x)$. 
    %Moreover 
    Recall that 
    \begin{align}
        \Phi = \left(JJ^{\top}\right)^{-1}. \notag
    \end{align}
    %Now, we compute 
    Then the entries of $I-\Pi$ are:
    \begin{align}
        [I - \Pi]_{lk} = \sum_{i,j} (\Phi^{-1})_{ij} (\nabla \xi_i)_{l} (\nabla \xi_j)_{k} = J^{\top}_{l}\Phi^{-1}J_{k}, %= \textsf{Tr}\left(J^{\top}_{l}\Phi^{-1}J_{k}\right) = \textsf{Tr}\left(\Phi^{-1}J_{k}J^{\top}_{l}\right), 
    \end{align}
    where $J_k$ and $ J_l$ are the $k$th and $l$th rows of $J$. Now, let $J = U_{r}\Sigma_{r}V^{\top}_{r}$ be the truncated SVD of $J$ where $r$ is the rank of $J$. Then $\Phi = \left(JJ^{\top}\right)^{-1} = U_{r}(\Sigma_{r}\Sigma_{r}^{\top})^{-1}U^{\top}_{r}$. Moreover, note that we can write each column of $J$ in terms of the columns of $U$ with coefficients from the rows of $V$. The $k$th column of $J$ is therefore:
    \begin{align}
        J_k = \sum_{\alpha=1}^{r}\sigma_{\alpha}v_{\alpha k}u_{\alpha}. 
    \end{align}
    Now we compute:
    \begin{align}
        \Phi^{-1}J_{k} &= U_{r}(\Sigma_{r}\Sigma_{r}^{\top})^{-1}U^{\top}_{r}\sum_{\alpha=1}^{r}\sigma_{\alpha}v_{\alpha k}u_{\alpha} \\
        &= \sum_{\alpha=1}^{r}\sigma_{\alpha}v_{\alpha k}U_{r}(\Sigma_{r}\Sigma_{r}^{\top})^{-1}U^{\top}_{r}u_{\alpha} \notag\\
        &= \sum_{\alpha=1}^{r}\sigma_{\alpha}v_{\alpha k}u_{\alpha}\sigma_{\alpha}^{-2} = \sum_{\alpha=1}^{r}\sigma_{\alpha}^{-1}v_{\alpha k}u_{\alpha}. \notag
    \end{align}
    Then $[I- \Pi]_{lk}$ is given by:  
    \begin{align}
 [I - \Pi]_{lk} &= J^{\top}_{l}\Phi^{-1}J_{k} = 
%        \textsf{Tr}
        \left(\sum_{\beta=1}^{r}\sigma_{\beta}v_{\beta l}u_{\beta}^{\top}\sum_{\alpha=1}^{r}\sigma_{\alpha}^{-1}v_{\alpha k}u_{\alpha}\right) \\
        &= \left(\sum_{1 \leq \alpha,\beta \leq r}\sigma_{\beta}v_{\beta l}\sigma_{\alpha}^{-1}v_{\alpha k}u_{\beta}^{\top}u_{\alpha}\right) = \left(\sum_{1 \leq \alpha \leq r}\sigma_{\alpha}v_{\alpha l}\sigma_{\alpha}^{-1}v_{\alpha k}\right) \notag\\
        &= \sum_{1 \leq \alpha \leq r}v_{\alpha l}v_{\alpha k} = \left[V_{r}V_{r}^{\top}\right]_{lk}. \notag
    \end{align}
    Thus, $[I - \Pi]_{lk} = \left[V_{r}V_{r}^{\top}\right]_{lk}$ so $I - \Pi = V_{r}V_{r}^{\top}$. 
    
    Now recall that \eqref{eq: orthogonality condition} states 
    \begin{equation*}
    D\xi(x)\nabla V_1(x) = J\nabla V_1 =  U_r\Sigma_rV_r^\top\nabla V_1 =0.
    \end{equation*}
    Since columns of $U_r\Sigma_r$ are linearly independent, the last equation is equivalent to
    \begin{equation*}
        V_r^\top\nabla V_1 = 0.
    \end{equation*}
    This, in turn, by the linear independence of columns of $V_r$, is equivalent to  $[I - \Pi] \nabla V_1 = V_rV_r^\top\nabla V_1 = 0.$
    Thus \eqref{eq: orthogonality condition} is equivalent to \eqref{eq: orthogonality condition 2}. 
    %By reversing the same argument we get that \eqref{eq: orthogonality condition 2} implies \eqref{eq: orthogonality condition}. 
\end{proof}

\end{document}
